\documentclass[11pt]{article}
\usepackage[utf8]{inputenc}
\usepackage{amsmath,amssymb,amsthm}
\usepackage[margin=1in]{geometry}
\usepackage{hyperref}
\usepackage{xcolor}

\newcommand{\tideref}[2][]{\href{https://therisensea.org/tides/#2/}{\texttt{#2}}}

\title{Tide retrospective: flat perturbations are invisible\\
(germbij Proposition 4.1)}
\author{Tide loop (Claude Fable 5, with GPT-5.6 Sol deliberation)}
\date{2026-08-10}

\begin{document}
\maketitle

\section{Setting}

The germbij note's identifiability programme proved the positive half
of its story: analytic potentials are pinned by their expectation
values, because an analytic pencil difference obeys a sector lower
bound $\kappa t^{\gamma}$ that contradicts superpolynomial decay. The
note's Proposition~4.1 is the negative half, and the reason the
analyticity hypothesis is not decoration: a \emph{flat} perturbation
--- smooth, nonnegative, vanishing to infinite order at the minimum
--- changes every observable integral by less than any power of $1/t$.
Without it, the identifiability theorem reads as a technical
convenience; with it, analyticity is exposed as the exact boundary of
what expectation values can see.

This was scoped as a candidate multi-tide programme, but scoping
collapsed it to a single tide: no expansion machinery is needed at
all. The pointwise identity
\[
e^{-tL} - e^{-t(L+f)} \;=\; e^{-tL}\,\bigl(1 - e^{-tf}\bigr),
\qquad 0 \le 1 - e^{-s} \le s,
\]
together with a \emph{global} polynomial domination of $f$, reduces
the whole difference to one closed-form Gaussian moment that the
seabed already owns.

\section{The thing formalised}

Two declarations in \texttt{Laplace/OneD/FlatInvisible.lean},
sorry-free with zero warnings. The main
theorem\footnote{\texttt{flat\_perturbation\_invisible}.}: if
$f$ is continuous with $0 \le f \le M$ and flat at $0$ (for every $n$
there are $C \ge 0$, $\delta > 0$ with $f(x) \le C x^{2n}$ on
$|x| \le \delta$), then for every continuous compactly supported
$\varphi$ and every $N$ there are $K \ge 0$ and $T \ge 1$ with
\[
\Bigl|\,\int \varphi\, e^{-t x^2/2}\,dx
 \;-\; \int \varphi\, e^{-t(x^2/2 + f)}\,dx \,\Bigr|
 \;\le\; \frac{K}{t^N}
\qquad (t \ge T).
\]
The corollary\footnote{\texttt{flat\_perturbation\_superpolynomial}.}
packages this as $\mathrm{IsLittleO}$: the difference is
$o(t^{-N})$ for every $N$ --- precisely the hypothesis shape that
\texttt{lower\_bound\_not\_superpolynomial} in
\texttt{Laplace/Decay.lean} refutes for analytic pencil differences.
The two halves of the note's story now meet in the same vocabulary:
analytic differences \emph{cannot} decay superpolynomially, flat ones
\emph{do}.

\section{Proof strategy}

Three moves, none requiring new infrastructure. First, a sup bound
$M_\varphi$ for $|\varphi|$ from compact support. Second, the global
domination: near $0$ flatness gives $f \le C x^{2N+2}$; away from
$0$ boundedness gives $f \le M \le (M/\delta^{2N+2})\,x^{2N+2}$, so
$f \le (C + M/\delta^{2N+2})\,x^{2N+2}$ everywhere. Third, the
pointwise bound
\[
\bigl|\varphi\,e^{-tx^2/2} - \varphi\,e^{-t(x^2/2+f)}\bigr|
 \;\le\; M_\varphi\, e^{-tx^2/2}\, t f
 \;\le\; M_\varphi (C{+}D)\, t\, x^{2N+2} e^{-tx^2/2},
\]
integrated by the seabed's closed-form moment
(\texttt{integral\_pow\_mul\_exp\_neg\_t\_sq\_half}), which yields
$c_N\, t \cdot t^{-(N+3/2)} \le c_N\, t^{-N}$ for $t \ge 1$. The
numerical check (run \emph{before} the consult, per the honesty rule)
confirmed superpolynomial decay for $f = e^{-1/x^2}$:
$|\Delta(t)|$ at $t = 20, 80, 320$ is
$8.4\times10^{-3}, 5.1\times10^{-5}, 4.2\times10^{-10}$, with
$\Delta\cdot t^3 \to 0$ and $\Delta \cdot t^5$ collapsing after onset.

\section{Roadblocks and resolutions}

One genuinely new failure mode. The double-factorial notation
$n\text{\textexclamdown\textexclamdown}$ needs \texttt{open scoped
Nat}, but that scope also defines $\varphi$ as notation for
\texttt{Nat.totient} --- so the theorem binder
$\{f\;\varphi : \mathbb{R} \to \mathbb{R}\}$ stopped parsing
(\emph{unexpected token, expected \}}) the moment the open was added.
Identifiers merely containing $\varphi$ survive; only the bare binder
dies. Resolution: skip the scoped open and write
\texttt{Nat.doubleFactorial} explicitly; the moment lemmas stated with
the notation still apply since it denotes the same constant. Promoted
to the repository error catalogue.

The only other friction was two known-catalogue items (the
\texttt{le\_or\_lt} rename, and a guessed integrability lemma name ---
the seabed's actual lemma is
\texttt{integrable\_pow\_mul\_exp\_neg\_mul\_sq} at $c = t/2$, bridged
by a one-line \texttt{ring} congruence). The $\mathrm{IsLittleO}$
corollary, including its rpow/npow bookkeeping, built on the first
attempt via \texttt{isLittleO\_iff\_tendsto'} and the ratio
$t^{-(N+1)}/t^{-N} = t^{-1}$.

\section{The deliberation}

GPT-5.6 Sol endorsed the candidate exactly as drafted: the even-power
flatness hypothesis (composes with the Gaussian moment lemmas without
rpow leakage), $T = 1$, the file name, and the all-orders $K/t^N$ form
as the quantitative core. Its one substantive addition was adopted:
\begin{quote}
The all-orders $O(t^{-N})$ statement faithfully captures ``no
asymptotic expansion coefficient changes.'' [\dots] I would make this
theorem the elementary quantitative result, and a short corollary
translating it into \texttt{Laplace.Decay}'s superpolynomial
vocabulary. That keeps the analytic proof independent of the
packaging.
\end{quote}
Its suggested helper decomposition was declined (each helper is used
once; the monolithic proof is $\sim$260 lines and readable), and its
recommendation to defer the classic witness $f = e^{-1/x^2}$ to a
separate declaration was accepted --- GPT supplied the flatness route
($s^n/n! \le e^s$ gives $e^{-1/x^2} \le n!\,x^{2n}$ globally, so
$C = n!$ and any $\delta$ works), which makes that follow-up tide
well-scoped if commissioned.

\section{Mathlib vs new}

Everything analytic was Mathlib or seabed lookup: the compact-support
sup bound, compact-support integrability, the inequality
$1 + s \le e^s$, the seabed's Gaussian moment and integrability
lemmas, and the tendsto characterisation of
$\mathrm{IsLittleO}$.\footnote{Respectively
\texttt{HasCompactSupport.}\allowbreak\texttt{exists\_bound\_of\_continuousOn},
\texttt{Continuous.}\allowbreak\texttt{integrable\_of\_hasCompactSupport},
\texttt{Real.add\_one\_le\_exp},
\texttt{integral\_pow\_mul\_exp\_neg\_t\_sq\_half} and
\texttt{integrable\_pow\_mul\_exp\_neg\_mul\_sq}, and
\texttt{isLittleO\_iff\_tendsto'}.} New content: the statement itself,
the global domination argument, and the assembly.

\section{What the next tide inherits}

The natural follow-up is the witness declaration: a concrete flat
$f = e^{-1/x^2}$ instance satisfying every hypothesis of
\texttt{flat\_perturbation\_invisible}, with flatness from the
factorial bound above (Mathlib's \texttt{expNegInvGlue} is the
reference point for continuity at $0$). With the witness in place the
note's Proposition~4.1 is fully witnessed rather than conditionally
quantified. Beyond that, the remaining germbij commissions are the
higher cumulant rungs of Theorem~3.1 and the multivariate \S7.4
programme.

\end{document}
